\scanpage{242}
et toute la situation de départ (62), en termes de \deleted{données} données de théorie des groupes plus explicites ou \textquotedblleft{}élémentaires\textquotedblright.

Quand on part de l'ensemble $S$, avec les deux sous-groupes

\[
\deleted{\unsure{G}}H\deleted{\unsure{0|Z|\sim}},Z\subset\unsure{\mathfrak S|\mathcal G}_S
\tag*{(65)}
\]

qui commutent l'un à l'autre, $Z$ commutatif, alors \deleted{la} \inserted{une} structure de \inserted{multi}graphe \inserted{$(S,\deleted{\unsure{A_0|A=}}(A_i)_{i\in\pi})$} sur $S$, admettant $H$ comme groupe d'automorphismes transitif \deleted{\unsure{sur les objets}} sur \inserted{les} $\overrightarrow{A_i}$ \inserted{$(i\in\pi)$}, et $Z$ comme groupe d'homothéties, est connue quand on connaît \inserted{($\pi$, et)} $A_{\bullet}(=\bigcup_{i\in\pi}A_i)\subset\mathfrak P_{\{1,2\}}(S)$,

\[
A\subset\underbrace{\mathfrak P_{\{1,2\}}}_{\substack{\text{parties à }1\\\text{ou }2\text{ élém.}}}(S),
\tag*{(66)}
\]

\deleted{qui doit être nécessairement un espace homogène} Car on récupère alors les $A_i$ $(i\in\pi)$ comme \deleted{\unsure{sous les groupes|sous le groupe}} les orbites de $H$ dans $A$, et on aura

\[
\deleted{(67)\qquad\widetilde H=H\mathbin{.} Z\subset\mathfrak S_S\quad -}
\]

\[
\pi\simeq A/H\quad .
\tag*{(67.)}
\]

\deleted{car on retrouve alors} \deleted{Pour \deleted{Explicitons} les conditions sur les données \inserted{Nous explicitons} (65) (66) pour qu'elles correspondent bien à un multigraphe \deleted{homogène} à groupe \deleted{d'automorphe} \inserted{d'homothéties} $Z$ et à groupe d'automorphismes homogène $A$, on trouve \deleted{que} \inserted{d'abord}}

\deleted{(67) $A$ stable par \unsure{$\sigma$} i.e. $H$ est un groupe d'autom pour $(S,A)$}

\[
\deleted{(68)\qquad\forall(s,t)\in S\times S,\ \inserted{\lambda\in Z,}\ \{s,t\}\in A\Longrightarrow\{\lambda s,t\}\in A}
\]

\deleted{--- condition qui est bien sûr valable pour $H$}

% archive sheet number: 241
